% ---------------------------------------------------------------
%  Muc II.3 -- Cac giai phap (TRONG TAM)
% ---------------------------------------------------------------
\subsection{Các giải pháp đã thực hiện}

Sáng kiến được thực hiện qua sáu giải pháp nối tiếp nhau, mỗi giải pháp giải
quyết một mắt xích của quy trình: từ cách \textit{đặt tên} câu hỏi, cách
\textit{viết} câu hỏi, cách \textit{quy định} ma trận, cách \textit{chọn} câu,
cách \textit{in} ra đề, đến cách \textit{tổ chức} cho học sinh làm bài.

\subsubsection{Giải pháp 1 --- Chuẩn hoá mã định danh câu hỏi}

Đây là giải pháp nền, quyết định toàn bộ những giải pháp còn lại. Nhận xét
xuất phát rất đơn giản: \textit{muốn máy tính chọn câu hỏi thay cho giáo viên
thì mỗi câu hỏi phải có một địa chỉ mà máy đọc được}. Vì vậy tôi quy định mọi
câu hỏi trong ngân hàng đều mang một mã định danh thống nhất:

\begin{figure}[H]
\centering
\begin{tikzpicture}[font=\small, node distance=0pt]
\node[font=\ttfamily\large] (id) at (0,0) {L10\_C1\_B2\_NB017\_MC\_A};
\node[font=\scriptsize, align=center] (a) at (-4.2,-1.6) {Khối lớp\\ (10/11/12)};
\node[font=\scriptsize, align=center] (b) at (-2.6,-2.4) {Chương};
\node[font=\scriptsize, align=center] (c) at (-1.2,-1.6) {Bài};
\node[font=\scriptsize, align=center] (d) at (0.4,-2.4) {Mức độ\\ NB/TH/VD/VDC};
\node[font=\scriptsize, align=center] (e) at (2.2,-1.6) {Mã năng lực\\ (lấy từ Curriculum)};
\node[font=\scriptsize, align=center] (f) at (3.9,-2.4) {Loại câu\\ MC/TF/SA/TL};
\node[font=\scriptsize, align=center] (g) at (5.2,-1.6) {Phiên bản};
\draw[gray] (-4.05,-0.25) -- (a);
\draw[gray] (-3.05,-0.25) -- (b);
\draw[gray] (-2.05,-0.25) -- (c);
\draw[gray] (-0.75,-0.25) -- (d);
\draw[gray] (0.55,-0.25) -- (e);
\draw[gray] (1.75,-0.25) -- (f);
\draw[gray] (2.75,-0.25) -- (g);
\end{tikzpicture}
\caption{Cấu trúc mã định danh của một câu hỏi trong ngân hàng}
\end{figure}

Bốn quy tắc kèm theo:

\begin{enumerate}[label=(\arabic*), leftmargin=1.2cm, itemsep=3pt]
  \item \textbf{Mã năng lực lấy nguyên văn từ bảng Curriculum}, không được tự
        đặt. Nhờ vậy mỗi câu hỏi luôn gắn với một yêu cầu cần đạt cụ thể của
        Chương trình 2018, và \textit{không thể vượt khung chương trình}: phạm
        vi kiến thức đã bị khoá ngay trong tên của câu hỏi.
  \item \textbf{Câu Đúng/Sai ra theo chương, không theo bài} --- mã rút gọn
        \texttt{L10\_C1\_TF\_A}. Theo quy ước chuyên môn của tôi, một câu
        Đúng/Sai gồm bốn ý ở bốn mức độ nhận thức khác nhau (nhận biết, thông
        hiểu, vận dụng, vận dụng cao), nên khi lên ma trận nó được tính là bốn
        câu chứ không phải một.
  \item \textbf{Mức vận dụng cao mượn mã năng lực của mức vận dụng.} Bảng
        Curriculum chỉ có ba mức nhận biết, thông hiểu, vận dụng; việc một câu
        được ra ở mức vận dụng hay vận dụng cao do bước chọn câu quyết định,
        không sinh thêm mã mới. Quy ước này giúp bảng Curriculum không phình
        ra vô ích.
  \item \textbf{Mã đã phát hành thì không sửa.} Muốn thay nội dung thì tạo
        phiên bản mới (\texttt{\_A} $\to$ \texttt{\_B} $\to$ \texttt{\_C}). Nhờ
        vậy một đề đã in ra từ hai năm trước vẫn tra ngược lại được đúng câu
        hỏi gốc.
\end{enumerate}

Toàn hệ thống --- cơ sở dữ liệu, giao diện, phần mềm điều phối --- đều làm việc
bằng mã định danh này chứ không bằng tên bài tiếng Việt, và \textbf{trí tuệ
nhân tạo tuyệt đối không được phép tự tạo hay tự suy luận ra mã}.

\subsubsection{Giải pháp 2 --- Viết ngân hàng câu hỏi bằng Python}

\textbf{a) Ba tầng dữ liệu}

Ngân hàng không phải một tệp câu hỏi mà là ba tầng dữ liệu nối nhau:

\begin{center}
\begin{tabular}{|p{2.6cm}|p{5.2cm}|p{5.6cm}|}
\hline
\textbf{Tầng} & \textbf{Nội dung} & \textbf{Ví dụ} \\ \hline
Phân phối chương trình & Danh sách bài của từng lớp kèm \textbf{số tiết} và mốc kiểm tra giữa kì, cuối kì &
\texttt{L10\_C1\_B2}, tiết ``23, 25--26'' $\to$ 3 tiết \\ \hline
Curriculum & Các yêu cầu cần đạt của từng bài, chia theo ba mức NB/TH/VD &
\texttt{L10\_C1\_B2\_NB017} \\ \hline
Mapping & Cầu nối: mỗi yêu cầu cần đạt có những hàm sinh câu hỏi nào, thuộc loại nào &
\texttt{L10\_C1\_B2\_NB017\_MC\_A} \\ \hline
\end{tabular}
\end{center}
\textit{Bảng 4. Ba tầng dữ liệu của ngân hàng đề}

\textbf{b) Câu hỏi là một hàm, không phải một đoạn văn bản}

Điểm khác biệt căn bản của sáng kiến nằm ở đây. Mỗi câu hỏi được viết thành một
hàm Python có tên trùng đúng mã định danh của nó. Hàm nhận vào số câu cần sinh
và trả về đoạn \LaTeX\ hoàn chỉnh gồm đề bài, các phương án, đánh dấu phương án
đúng và lời giải chi tiết:

\begin{lstlisting}
def L10_C1_B2_NB017_MC_A_01(socau, socot):
    ket_qua = ""
    for i in range(socau):
        a = random.randint(2, 9)          # sinh so lieu ngau nhien
        b = random.choice([x for x in range(2, 12) if x % a == 0])
        dap_an_dung = b // a              # dap an do PYTHON tinh ra
        ...                               # dung 3 phuong an nhieu
        ket_qua += cau_hoi_latex
    return ket_qua
\end{lstlisting}

Ba hệ quả quan trọng:

\begin{itemize}[leftmargin=1.2cm, itemsep=3pt]
  \item \textbf{Đáp án không thể sai.} Đáp án và lời giải không phải do ai gõ
        vào hay do trí tuệ nhân tạo đoán ra, mà do chính Python tính bằng công
        thức toán học từ bộ số liệu vừa sinh. Mỗi lần chạy là một lần tính
        lại, nên số liệu đổi thì đáp án đổi theo, luôn khớp.
  \item \textbf{Số liệu luôn ``đẹp''.} Phần công sức chuyên môn của giáo viên
        nằm ở chỗ ràng buộc miền giá trị của tham số sao cho mọi biến thể đều
        có nghiệm đẹp, đúng tầm học sinh và không rơi vào trường hợp suy biến.
        Máy không làm thay được việc này.
  \item \textbf{Một mã, nhiều cách ra đề.} Cùng một mã định danh có thể có
        nhiều hàm khác nhau, đánh số \texttt{\_01}, \texttt{\_02},
        \texttt{\_03}\dots\ --- mỗi hàm là một bối cảnh ra đề khác. Khi sinh
        đề, hệ thống chọn ngẫu nhiên một trong các hàm đó. Hậu tố này chỉ tồn
        tại bên trong tệp Python, không lộ ra bất kì nơi nào khác, nên việc bổ
        sung độ phong phú cho ngân hàng không làm thay đổi mã định danh đã
        phát hành.
\end{itemize}

\textbf{c) Kiểm thử tự động toàn bộ ngân hàng}

Ngân hàng càng lớn thì càng khó rà bằng mắt. Tôi viết thêm bộ kiểm thử tự động
gọi lần lượt toàn bộ các hàm sinh câu hỏi và kiểm tra kết quả trả về có đúng
định dạng, có đủ đáp án hay không. Bộ kiểm thử này đã phát hiện một lỗi thật mà
đọc bằng mắt không thấy: một hàm kết thúc bằng \texttt{return} mà quên tên biến
nên trả về rỗng, khiến câu hỏi biến mất khỏi đề trong suốt một thời gian dài.

\chenanh[0.92]{03-ham-sinh-cau-hoi.png}{Một hàm sinh câu hỏi trong ngân hàng Python, kèm phần tính đáp án và lời giải}

\subsubsection{Giải pháp 3 --- Đưa ma trận đề vào bảng cấu hình}

Ma trận đề không được viết lẫn vào mã chương trình mà đặt riêng trong một bảng
cấu hình, để khi Bộ hoặc Sở thay đổi cấu trúc đề thì chỉ cần sửa số trong bảng.
Bảng quy định cho từng loại bài kiểm tra (hệ số 1 --- kiểm tra thường xuyên,
hệ số 2 và 3 --- kiểm tra định kì) số lượng câu của mỗi phần và tỉ lệ bốn mức
độ nhận thức:

\begin{lstlisting}
"HeSo2_HeSo3": {
  "trac_nghiem":     {"so_luong": 12, "ty_le_muc_do": {"NB":0.4,"TH":0.3,"VD":0.2,"VDC":0.1}},
  "dung_sai_cau_lon":{"so_luong":  2, "ty_le_muc_do": {"NB":0.25,"TH":0.25,"VD":0.25,"VDC":0.25}},
  "tra_loi_ngan":    {"so_luong":  4, "ty_le_muc_do": {"NB":0,"TH":0,"VD":0.5,"VDC":0.5}},
  "tu_luan":         {"so_luong":  3, "ty_le_muc_do": {"NB":0,"TH":0,"VD":0.7,"VDC":0.3}}
}
\end{lstlisting}

Một bảng thứ hai quy định thang điểm 10 chia theo phần: trắc nghiệm nhiều lựa
chọn 3 điểm, đúng/sai 2 điểm, trả lời ngắn 2 điểm, tự luận 3 điểm; trong mỗi
phần thì chia đều cho từng câu. Cách tách này cho phép thay đổi trọng số các
phần mà không phải sửa một dòng mã nào.

\ghichu{Nếu tổ chuyên môn của chị thống nhất một ma trận khác (ví dụ hệ số 1
chỉ 6 câu trắc nghiệm), sửa số trong bảng rồi chụp lại ảnh cho khớp với bài
viết.}

\chenanh[0.92]{04-bang-ma-tran.png}{Bảng cấu hình ma trận đề và bảng thang điểm --- thay đổi cấu trúc đề chỉ cần sửa số ở đây}

\subsubsection{Giải pháp 4 --- Thuật toán dựng ma trận chi tiết và chọn câu hỏi}

Đây là phần lõi và cũng là phần được chỉnh sửa nhiều lần nhất. Từ một yêu cầu
của người dùng, hệ thống đi qua sáu bước, \textbf{tất cả đều là chương trình
máy tính, không có trí tuệ nhân tạo tham gia}:

\textbf{Bước 1 --- Xác định phạm vi.} Đọc phân phối chương trình theo lớp.
Nếu là kiểm tra thường xuyên theo chương thì lấy toàn bộ bài của chương đó;
nếu là giữa kì hoặc cuối kì thì lấy từ đầu học kì đến đúng bài có mốc kiểm tra
tương ứng. Bước này trả về danh sách bài kèm \textbf{số tiết} của từng bài.

\textbf{Bước 2 --- Chọn bài cho câu Đúng/Sai trước tiên.} Vì mỗi câu Đúng/Sai
chiếm trọn bốn mức độ nên phải xếp chỗ cho nó trước, nếu không sẽ vỡ ma trận.
Các bài trong phạm vi được sắp theo số tiết giảm dần, lấy từ trên xuống; bốn ý
của một câu Đúng/Sai đều nằm trong cùng một bài.

\textbf{Bước 3 --- Trừ chỉ tiêu ngay tại bài đã chọn.} Một câu Đúng/Sai đã
chiếm 1 nhận biết + 1 thông hiểu + 1 vận dụng + 1 vận dụng cao, nên phải trừ
đúng bốn ý đó ra khỏi chỉ tiêu của chính bài đó --- \textit{``đơn vị kiến thức
đó chia ra được 3 câu thì phải tính 1 ở đúng sai, chỉ còn 2''}. Trừ ở cấp bài
đã tự động làm tổng giảm đúng bằng số câu Đúng/Sai, tuyệt đối không trừ thêm
lần thứ hai ở cấp tổng.

\textbf{Bước 4 --- Chia câu về từng bài theo tỉ lệ số tiết.} Các câu mức nhận
biết và thông hiểu được chia về từng bài \textit{theo tỉ lệ số tiết của bài đó
trong phân phối chương trình} --- bài dạy nhiều tiết thì nhiều câu, đúng
nguyên tắc trọng số của ma trận đề. Phần dư sau khi làm tròn xuống được rải
lần lượt cho bài nhiều tiết nhất trước.

\textbf{Bước 5 --- Rải mức vận dụng và vận dụng cao.} Chọn vận dụng cao trước,
\textbf{mỗi bài nhiều nhất một câu vận dụng cao} để không dồn hết câu khó vào
một đơn vị kiến thức. Sau đó mới rải câu vận dụng, ưu tiên những bài chưa có
câu vận dụng cao. Tỉ lệ vận dụng so với vận dụng cao lấy đúng từ bảng cấu hình
ở Giải pháp 3, không ép cứng trong mã.

\textbf{Bước 6 --- Chọn mã câu hỏi cuối cùng.} Đến đây ma trận đã nói rõ ``cần
một câu mức thông hiểu, thuộc yêu cầu cần đạt số 014 của bài 1, loại trắc
nghiệm nhiều lựa chọn''. Bước chọn câu tra bảng Mapping để tìm mã hàm sinh
tương ứng, bỏ qua những mã đã dùng trong đề hiện tại, và xoay vòng giữa các
phiên bản A, B, C để đề nào cũng khác đề nào.

Sau đó câu hỏi sinh ra còn qua một bước kiểm tra tự động: đúng chương, đúng
bài, đúng mức độ, đúng loại câu, không trùng câu trong cùng một đề.

\textbf{Chế độ nháp.} Nếu ngân hàng chưa đủ câu cho một ô nào đó của ma trận,
hệ thống không dừng cả đề mà chèn dòng ``[THIẾU CÂU HỎI]'' vào đúng vị trí.
Đây là lựa chọn có chủ ý: dòng đó vừa cho giáo viên một đề dùng tạm được, vừa
là chỉ dấu rõ ràng cho biết cần bổ sung hàm sinh câu hỏi nào tiếp theo.

\chenanh[0.92]{05-ma-tran-de-sinh-ra.png}{Ma trận chi tiết do hệ thống dựng ra cho một đề kiểm tra định kì}

\clearpage
\subsubsection{Giải pháp 5 --- Sinh đề, lời giải và nhiều mã đề bằng \LaTeX}

\textbf{a) Đề và lời giải từ cùng một nguồn}

Các câu hỏi đã chọn được ghép thành một tệp \LaTeX\ theo mẫu đề thi, rồi biên
dịch ra PDF. Bản lời giải không được sinh lại lần nữa mà lấy chính tệp
\LaTeX\ của đề vừa in, chỉ đổi chế độ hiển thị từ ``đề thi'' sang ``lời giải''.
Nhờ vậy \textbf{lời giải chắc chắn ứng đúng với đề học sinh đang cầm trên tay},
không bao giờ xảy ra tình trạng lệch câu giữa đề và đáp án --- một lỗi rất
thường gặp khi làm thủ công.

Đáp án cũng được trích ngay tại thời điểm sinh đề, đọc trực tiếp từ các dấu
hiệu trong mã \LaTeX\ (phương án được đánh dấu đúng, đáp số của câu trả lời
ngắn, phần lời giải), rồi lưu lại thành một tệp đáp án đi kèm. Khi chấm bài,
hệ thống so với tệp đáp án này, nên kết quả chấm khớp 100\% với đề đã phát.

\textbf{b) Nhiều mã đề cùng cấu trúc, chỉ khác số liệu}

Đây là chức năng được giáo viên đánh giá cao nhất khi dùng thử. Vì mỗi câu hỏi
là một hàm sinh, nên chỉ cần gọi lại hàm đó là có ngay một câu mới:
\textbf{cùng bài, cùng yêu cầu cần đạt, cùng mức độ nhận thức, cùng dạng toán
--- chỉ khác số liệu}. Từ đó hệ thống in ra bao nhiêu mã đề cũng được, và các
mã đề ấy tương đương nhau tuyệt đối về nội dung lẫn độ khó.

Ý nghĩa thực tế với giáo viên: \textit{lấy ra là dùng kiểm tra được ngay, không
phải ngồi tự thay số để tạo đề mới, cũng không phải giải lại để kiểm tra đáp
án}. Ý nghĩa với học sinh: hai em ngồi cạnh nhau nhận hai đề khác nhau thật
sự, nhưng không em nào thiệt hơn em nào.

Cũng từ cơ chế đó, trang kết quả sau khi làm bài có thêm nút ``Làm đề khác cùng
cấu trúc''. Học sinh bấm một cái là có ngay một đề mới y hệt về cấu trúc để
luyện lại phần mình vừa làm sai --- và đường đi này \textbf{không gọi trí tuệ
nhân tạo}, vì tham số của đề cũ đã được lưu sẵn, nên vừa nhanh vừa không tốn
chi phí.

\chenanh[0.92]{06-hai-ma-de.png}{Hai mã đề của cùng một cấu trúc: cùng dạng toán, cùng mức độ, chỉ khác số liệu}

\chenanh[0.92]{07-de-va-loi-giai-pdf.png}{Tệp PDF đề thi và tệp PDF lời giải tương ứng do hệ thống sinh ra}

\subsubsection{Giải pháp 6 --- Tổ chức làm bài, chấm điểm và theo dõi kết quả}

\textbf{a) Làm bài trực tuyến và chấm ngay.} Học sinh mở đề trên web, làm các
phần trắc nghiệm nhiều lựa chọn, đúng/sai và trả lời ngắn ngay trên trình
duyệt. Bấm nộp là có điểm ngay cùng với đáp án đúng của từng câu. Việc chấm do
chương trình đối chiếu với tệp đáp án đã lưu, \textbf{không gọi trí tuệ nhân
tạo}, nên vừa chính xác vừa không phát sinh chi phí dù bao nhiêu học sinh cùng
làm bài.

\textbf{b) Thang điểm 10 chia theo phần.} Điểm được tính theo đúng bảng trọng
số ở Giải pháp 3: trắc nghiệm nhiều lựa chọn 3 điểm, đúng/sai 2 điểm, trả lời
ngắn 2 điểm, tự luận 3 điểm, mỗi phần chia đều cho các câu trong phần đó. Hiện
tại phần tự luận 3 điểm học sinh làm ra giấy nộp cho giáo viên chấm, nên màn
hình kết quả ghi rõ điểm tự động là 7 điểm và chức năng chấm tự luận đang được
xây dựng --- ghi rõ như vậy để học sinh không hiểu nhầm là mình bị mất điểm.

\textbf{c) Giao đề và trả điểm qua Google Classroom.} Sau khi giáo viên kết
nối một lần tài khoản Google Classroom của mình, hệ thống tự đăng đề, lời giải
và điểm của từng học sinh về đúng lớp học thật trên Classroom. Học sinh mở
Classroom quen thuộc là thấy bài của mình, không phải học thêm một phần mềm
mới.

\textbf{d) Trang thống kê cho giáo viên.} Giáo viên xem được điểm trung bình
của lớp, danh sách điểm từng học sinh và --- quan trọng hơn cả --- những
chương, bài mà học sinh sai nhiều nhất. Đây chính là thông tin để điều chỉnh
việc dạy: bài nào cả lớp sai nhiều thì cần dạy lại, em nào sai lệch hẳn so với
lớp thì cần kèm riêng.

\chenanh[0.92]{15-lam-bai-nop-bai.png}{Màn hình làm bài trực tuyến và màn hình kết quả chấm ngay sau khi nộp}

\subsubsection{Công cụ và hạ tầng đã sử dụng}

Sáng kiến này là một phần mềm chạy thật trên Internet, nên tôi xin nêu rõ toàn
bộ công cụ đã dùng để đồng nghiệp muốn làm theo biết cần chuẩn bị những gì.

\begin{center}
\begin{longtable}{|p{3.6cm}|p{4.6cm}|p{6.2cm}|}
\hline
\textbf{Thành phần} & \textbf{Công cụ đã dùng} & \textbf{Dùng để làm gì} \\ \hline
\endhead
Ngôn ngữ lập trình & Python 3 & Viết hàm sinh câu hỏi và toàn bộ phần xử lí \\ \hline
Máy chủ web & FastAPI + Jinja2 & Cung cấp giao diện và các đường dẫn của trang web \\ \hline
Cơ sở dữ liệu, đăng nhập & Supabase (PostgreSQL) & Lưu tài khoản, lịch sử hội thoại, đề đã sinh, bài làm, điểm \\ \hline
Điều phối luồng & n8n (bản tự cài trên máy chủ riêng) & Nối các bước với nhau và là nơi duy nhất gọi tới trí tuệ nhân tạo \\ \hline
Mô hình ngôn ngữ & Gọi qua OpenRouter (Qwen, Gemini) & Chỉ để hiểu yêu cầu tiếng Việt của người dùng \\ \hline
Trình bày đề & \LaTeX, gói \texttt{ex\_test}, biên dịch bằng \texttt{pdflatex} & Sinh tệp PDF đề thi và lời giải đúng chuẩn văn bản đề \\ \hline
Máy chủ & Máy ảo (VPS) chạy Ubuntu, nginx, systemd & Nơi phần mềm chạy 24/24 \\ \hline
Tên miền & \texttt{nganhangdechv.tech} (tự mua) & Địa chỉ để giáo viên và học sinh truy cập \\ \hline
Tích hợp ngoài & Google Classroom API, Google Drive API & Giao đề, đáp án và trả điểm về lớp học thật \\ \hline
Quản lí mã nguồn & Git và GitHub & Lưu lịch sử sửa đổi, đồng bộ giữa máy cá nhân và máy chủ \\ \hline
Soạn thảo & Visual Studio Code & Nơi viết toàn bộ mã nguồn \\ \hline
\end{longtable}
\end{center}
\textit{Bảng 5. Công cụ và hạ tầng đã sử dụng}

Trong đó, hai thành phần cần giải thích thêm:

\textbf{n8n} là một phần mềm nối các bước xử lí thành một luồng, thao tác bằng
cách kéo thả các khối trên màn hình thay vì viết mã. Tôi cài n8n trên máy chủ
của mình (bản tự cài, không phải bản thuê theo tháng). Trong hệ thống này, n8n
giữ đúng một vai trò: nhận yêu cầu tiếng Việt, chuyển cho mô hình ngôn ngữ hiểu,
rồi gọi sang phần xử lí bằng Python. Mọi việc có quy tắc rõ ràng đều đã được
chuyển hết về Python, đúng nguyên tắc ``ưu tiên mã nguồn hơn trí tuệ nhân tạo''.

\textbf{Máy ảo (VPS)} là một máy tính thuê trên Internet, chạy liên tục, để
phần mềm luôn sẵn sàng khi học sinh truy cập bất kể giờ nào. Cùng với tên miền,
đây là hai khoản chi phí duy nhất của sáng kiến, do tôi tự chi trả.

\ghichu{Nếu muốn ghi rõ chi phí, điền số tiền thật mỗi năm cho máy ảo và tên
miền vào đoạn trên --- con số cụ thể sẽ làm phần ``hiệu quả kinh tế'' ở Mục
II.5 thuyết phục hơn nhiều.}

\textbf{Mã nguồn công khai.} Toàn bộ mã nguồn của hệ thống, kèm hơn 20 tài liệu kĩ
thuật mô tả kiến trúc, chuẩn mã định danh, thuật toán dựng ma trận và nhật kí
toàn bộ các lần cải tiến, được công khai tại:

\begin{center}
\setlength{\fboxsep}{8pt}
\fbox{\begin{minipage}{0.85\linewidth}\centering
\textbf{\url{https://github.com/lantuan/NganHangDe}}
\end{minipage}}
\end{center}

Tôi công khai mã nguồn vì hai lí do. Thứ nhất, để hội đồng và đồng nghiệp kiểm
chứng được rằng sản phẩm là có thật và do tôi làm, chứ không chỉ là mô tả trên
giấy --- riêng nhật kí sửa đổi trong mã nguồn đã ghi lại từng lần cải tiến ở
Bảng 6 dưới đây. Thứ hai, để giáo viên nào muốn làm theo có thể lấy về dùng
lại, không phải bắt đầu từ con số không.

\subsubsection{Quá trình cải tiến}

Sáng kiến không hoàn chỉnh ngay từ bản đầu tiên. Toàn bộ quá trình được ghi lại
trong sổ tay kĩ thuật của hệ thống theo từng phiên bản; dưới đây là những lần
cải tiến đáng kể nhất, đi kèm lí do --- vì chính những lần sửa này mới là nội
dung có giá trị chia sẻ với đồng nghiệp.

\begin{center}
\begin{longtable}{|p{1.4cm}|p{5.6cm}|p{6.4cm}|}
\hline
\textbf{Lần} & \textbf{Vấn đề gặp phải} & \textbf{Cách khắc phục} \\ \hline
\endhead
1 & Bản thiết kế ban đầu định dùng nhiều trí tuệ nhân tạo cho từng khâu (phân
tích yêu cầu, dựng ma trận, chọn câu) --- kết quả không ổn định, mỗi lần chạy
một khác &
Rút xuống chỉ còn ba chỗ dùng trí tuệ nhân tạo, toàn bộ khâu còn lại chuyển
sang chương trình máy tính. Đây là lần cải tiến quan trọng nhất, đặt ra nguyên
tắc xuyên suốt của hệ thống \\ \hline

2 & Đáp án phải nhập tay riêng nên hay lệch với đề &
Phát hiện đáp án đã nằm sẵn trong mã \LaTeX\ của câu hỏi; viết chương trình tự
trích ra ngay lúc sinh đề. Việc trích không dùng được cách tìm kiếm thông
thường vì công thức toán có ngoặc lồng nhau, phải viết bộ đếm ngoặc riêng \\ \hline

3 & Lời giải bị lọt vào đề của học sinh &
Bổ sung mốc chặn khi cắt phần đề bài, đồng thời lược bỏ các lệnh \LaTeX\ mà
trình duyệt không hiển thị được \\ \hline

4 & Mọi câu Đúng/Sai in ra PDF bị dính liền thành một đoạn văn &
Truy ra một lệnh \LaTeX\ được sinh ra nhưng không hề tồn tại trong gói lệnh
đang dùng; sửa ở thư viện sinh câu hỏi để luôn dùng lệnh đúng \\ \hline

5 & Toàn bộ câu trả lời ngắn bị sinh nhầm thành câu trắc nghiệm ở mọi đề đã
từng tạo &
Nguyên nhân là chỗ gọi hàm luôn truyền cùng một giá trị cho tham số dạng câu;
sửa lại nhánh gọi và bổ sung giá trị mặc định cho 29 hàm sinh \\ \hline

6 & Một hàm sinh câu hỏi trả về rỗng, câu hỏi biến mất khỏi đề mà không ai biết &
Viết bộ kiểm thử tự động chạy qua toàn bộ các hàm trong ngân hàng; lỗi bị phát
hiện ngay và từ đó mọi lần bổ sung câu hỏi mới đều phải qua bộ kiểm thử này \\ \hline

7 & Số câu do người dùng yêu cầu bị bỏ mất trên đường truyền, đề luôn ra theo
số mặc định &
Rà lại toàn bộ đường đi của dữ liệu, phát hiện một bước trung gian gán cứng giá
trị rỗng đè lên yêu cầu thật \\ \hline

8 & Trí tuệ nhân tạo hiểu ``chương 1'' thành ``lớp 1'' &
Chặn sớm ngay đầu vào: lớp không thuộc 10, 11, 12 thì báo lỗi rõ ràng; sau đó
đổi hẳn thành hỏi lại người dùng thay vì tự đoán \\ \hline

9 & Câu ``cho tôi thêm đề nữa'' bị hiểu là ngoài phạm vi &
Thay vì chỉnh trí tuệ nhân tạo cho hiểu, làm hẳn một đường đi không qua trí tuệ
nhân tạo: nút ``Làm đề khác cùng cấu trúc'' dùng lại tham số của đề cũ \\ \hline

10 & Số câu chia đều cho các bài, không đúng trọng số &
Chuyển sang chia theo \textbf{tỉ lệ số tiết} của từng bài trong phân phối
chương trình \\ \hline

11 & Câu Đúng/Sai bị trừ chỉ tiêu hai lần, làm hụt câu ở các phần khác &
Quy định trừ đúng một lần và trừ ngay tại bài mà câu Đúng/Sai đã chọn \\ \hline

12 & Các câu vận dụng cao dồn hết vào một đơn vị kiến thức &
Quy định mỗi bài nhiều nhất một câu vận dụng cao, chọn vận dụng cao trước rồi
mới rải vận dụng \\ \hline

13 & Học sinh mất phiên đăng nhập giữa chừng &
Bổ sung cơ chế tự làm mới phiên khi phiên hết hạn \\ \hline

14 & Đề và đáp án không hiện trong tài khoản Classroom của học sinh &
Ba nguyên nhân nối nhau: giao diện nuốt mất thông báo lỗi; chưa bật một dịch vụ
cần thiết của Google; và xin nhầm quyền (quyền đăng \textit{bài tập} khác quyền
đăng \textit{tài liệu}). Sửa lần lượt cả ba, đồng thời cho hiện lí do thất bại
ngay trên màn hình thay vì im lặng \\ \hline
\end{longtable}
\end{center}
\textit{Bảng 6. Những lần cải tiến đáng kể trong quá trình thực hiện}

Nhìn lại, hầu hết các lỗi đều thuộc một trong hai nhóm: \textit{giao cho trí
tuệ nhân tạo việc lẽ ra phải do chương trình làm}, và \textit{không có cơ chế
tự kiểm tra nên lỗi âm thầm tồn tại rất lâu}. Hai bài học rút ra tương ứng là
nguyên tắc phân vai ở Giải pháp 1 và bộ kiểm thử tự động ở Giải pháp 2.

\subsubsection{Vai trò của trợ lí trí tuệ nhân tạo trong quá trình thực hiện}

Tôi xin nói rõ và không né tránh: trong quá trình xây dựng hệ thống này, tôi có
sử dụng trợ lí trí tuệ nhân tạo (Claude của Anthropic) như một cộng sự kĩ
thuật. Cụ thể, trợ lí giúp tôi viết và sửa mã nguồn theo yêu cầu, dò tìm
nguyên nhân của các lỗi nêu ở bảng trên, và ghi lại sổ tay kĩ thuật của dự án.

Tuy nhiên, phần thuộc về chuyên môn sư phạm vẫn hoàn toàn do tôi quyết định và
chịu trách nhiệm: quy ước mã định danh câu hỏi, quy ước một câu Đúng/Sai gồm
bốn ý ở bốn mức độ, quy tắc trừ chỉ tiêu tại chính bài đã chọn, nguyên tắc chia
câu theo tỉ lệ số tiết, quy định mỗi bài nhiều nhất một câu vận dụng cao, ma
trận và thang điểm của từng loại bài kiểm tra, cùng toàn bộ nội dung toán học
và ràng buộc số liệu của từng hàm sinh câu hỏi. Nhiều lần trợ lí đề xuất cách
làm khác, tôi phải yêu cầu làm lại cho đúng quy ước sư phạm mà tôi đã đặt ra.

Chính trải nghiệm đó củng cố thêm quan điểm nêu ở Mục II.1: trí tuệ nhân tạo là
một công cụ rất mạnh để hiện thực hoá ý tưởng của giáo viên, nhưng không thay
được người giáo viên trong việc quyết định \textit{dạy gì, kiểm tra gì và kiểm
tra đến mức nào}. Với các đồng nghiệp muốn làm theo hướng này, tôi cho rằng
điều kiện cần không phải là biết lập trình giỏi, mà là \textbf{có quy ước
chuyên môn rõ ràng trước khi ngồi vào máy} --- vì đó chính là thứ mà công cụ
không tự nghĩ ra được.

\clearpage
